% GENERATED FILE. Edit canonical lessons, then run npm run generate:books.
% canonical-source-hash: fnv1a64:09472140f0a5fed1
% canonical-lessons: HI-C407-rajmarg, HI-C407-tol, HI-C407-naka, HI-C407-qafila, HI-C407-bahana

\chapter{On the Highway}
\label{ch:hi-on-the-highway}
\hlchaptermodality{407}

\begin{chapteropening}
\textbf{By the end of this chapter:} \emph{I can talk about a highway, its tolls and checkpoints, a convoy of lorries and a road washed away.}

Complete the last lesson of chapter 407: I can talk about a highway, its tolls and checkpoints, a convoy of lorries and a road washed away.
\end{chapteropening}

\section[\texorpdfstring{\hi{राजमार्ग} (rājmārg)}{राजमार्ग (rājmārg)}]{\hi{राजमार्ग} (rājmārg) — a highway}

\label{lesson:HI-C407-rajmarg}



\begin{quote}
Before the new one: say the Hindi for a helmet, then the Hindi for a licence.
\end{quote}

\subsection*{You'll want to know: \hi{राजमार्ग}}
\textbf{\hi{राजमार्ग}} — \emph{rājmārg} — \textquotedblleft{}a highway\textquotedblright{}.

\begin{quote}
\hi{यह} \hi{राजमार्ग} \hi{पहाड़ों} \hi{से} \hi{होकर} \hi{जाता} \hi{है।}
\end{quote}

\begin{quote}
\emph{yah rājmārg pahāṛoṁ se hokar jātā hai.}
\end{quote}

\begin{quote}
This highway goes through the mountains.
\end{quote}

Masculine.

\begin{cousinweb}
\textbf{\hi{राजमार्ग}} (\emph{rājmārg}) is a tatsama, Sanskrit \emph{rājamārga}, \textquotedblleft{}the king's road, a main road\textquotedblright{} (Monier-Williams), from \emph{rāja}, \textquotedblleft{}king\textquotedblright{}, and \emph{mārga}, \textquotedblleft{}road\textquotedblright{}; the national highways carry it on their signs.
\end{cousinweb}

\begin{grammarlens}[title={What you've built}]
One more word for getting around.
\end{grammarlens}

\subsection*{Your turn}
\begin{itemize}
\raggedright
  \item \emph{Say it:} \emph{rājmārg}
  \item \emph{Say it:} \emph{rājmārg}, once more
  \item \emph{Say it:} \emph{yah rājmārg pahāṛoṁ se hokar jātā hai.} — This highway goes through the mountains.
  \item \emph{Recall:} say the Hindi for a helmet, then the Hindi for a licence, then say \emph{rājmārg} again
\end{itemize}

\begin{tcolorbox}[breakable,colback=teal!4,colframe=teal!35!black,title={Before you move on}]
What does \hi{राजमार्ग} mean? (\textquotedblleft{}A highway\textquotedblright{}.) Say it once more.
\end{tcolorbox}
\section[\texorpdfstring{\hi{टोल} (ṭol)}{टोल (ṭol)}]{\hi{टोल} (ṭol) — a toll (on a road)}

\label{lesson:HI-C407-tol}



\begin{quote}
Before the new one: say the Hindi for a licence, then the Hindi for a highway.
\end{quote}

\subsection*{You'll want to know: \hi{टोल}}
\textbf{\hi{टोल}} — \emph{ṭol} — \textquotedblleft{}a toll (on a road)\textquotedblright{}.

\begin{quote}
\hi{राजमार्ग} \hi{पर} \hi{हमें} \hi{दो} \hi{जगह} \hi{टोल} \hi{देना} \hi{पड़ा।}
\end{quote}

\begin{quote}
\emph{rājmārg par hamẽ do jagah ṭol denā paṛā.}
\end{quote}

\begin{quote}
We had to pay a toll at two places on the highway.
\end{quote}

Masculine.

\begin{cousinweb}
\textbf{\hi{टोल}} (\emph{ṭol}) is borrowed from English \emph{toll}, the charge paid at a barrier on a highway or a bridge.
\end{cousinweb}

\begin{grammarlens}[title={What you've built}]
One more word for getting around.
\end{grammarlens}

\subsection*{Your turn}
\begin{itemize}
\raggedright
  \item \emph{Say it:} \emph{ṭol}
  \item \emph{Say it:} \emph{ṭol}, once more
  \item \emph{Say it:} \emph{rājmārg par hamẽ do jagah ṭol denā paṛā.} — We had to pay a toll at two places on the highway.
  \item \emph{Recall:} say the Hindi for a licence, then the Hindi for a highway, then say \emph{ṭol} again
\end{itemize}

\begin{tcolorbox}[breakable,colback=teal!4,colframe=teal!35!black,title={Before you move on}]
What does \hi{टोल} mean? (\textquotedblleft{}A toll (on a road)\textquotedblright{}.) Say it once more.
\end{tcolorbox}
\section[\texorpdfstring{\hi{नाका} (nākā)}{नाका (nākā)}]{\hi{नाका} (nākā) — a checkpoint (on a road)}

\label{lesson:HI-C407-naka}



\begin{quote}
Before the new one: say the Hindi for a highway, then the Hindi for a toll.
\end{quote}

\subsection*{You'll want to know: \hi{नाका}}
\textbf{\hi{नाका}} — \emph{nākā} — \textquotedblleft{}a checkpoint (on a road)\textquotedblright{}.

\begin{quote}
\hi{शहर} \hi{के} \hi{नाके} \hi{पर} \hi{पुलिस} \hi{गाड़ियाँ} \hi{रोक} \hi{रही} \hi{थी।}
\end{quote}

\begin{quote}
\emph{śahar ke nāke par pulis gāṛiyā̃ rok rahī thī.}
\end{quote}

\begin{quote}
At the checkpoint into the town the police were stopping cars.
\end{quote}

Masculine.

\begin{cousinweb}
\textbf{\hi{नाका}} (\emph{nākā}) goes back to \emph{*nākka}, \textquotedblleft{}like a nose, projecting\textquotedblright{}, a form CDIAL reconstructs (CDIAL 7037), which lists Hindi \emph{nākā}, \textquotedblleft{}the end of a road, a passage, a gap\textquotedblright{}: the narrow point where a road enters a town, and so the post that watches it.
\end{cousinweb}

\begin{grammarlens}[title={What you've built}]
One more word for getting around.
\end{grammarlens}

\subsection*{Your turn}
\begin{itemize}
\raggedright
  \item \emph{Say it:} \emph{nākā}
  \item \emph{Say it:} \emph{nākā}, once more
  \item \emph{Say it:} \emph{śahar ke nāke par pulis gāṛiyā̃ rok rahī thī.} — At the checkpoint into the town the police were stopping cars.
  \item \emph{Recall:} say the Hindi for a highway, then the Hindi for a toll, then say \emph{nākā} again
\end{itemize}

\begin{tcolorbox}[breakable,colback=teal!4,colframe=teal!35!black,title={Before you move on}]
What does \hi{नाका} mean? (\textquotedblleft{}A checkpoint (on a road)\textquotedblright{}.) Say it once more.
\end{tcolorbox}
\section[\texorpdfstring{\hi{क़ाफ़िला} (qāfilā)}{क़ाफ़िला (qāfilā)}]{\hi{क़ाफ़िला} (qāfilā) — a convoy; a caravan}

\label{lesson:HI-C407-qafila}



\begin{quote}
Before the new one: say the Hindi for a toll, then the Hindi for a checkpoint.
\end{quote}

\subsection*{You'll want to know: \hi{क़ाफ़िला}}
\textbf{\hi{क़ाफ़िला}} — \emph{qāfilā} — \textquotedblleft{}a convoy; a caravan\textquotedblright{}.

\begin{quote}
\hi{राजमार्ग} \hi{पर} \hi{ट्रकों} \hi{का} \hi{एक} \hi{लंबा} \hi{क़ाफ़िला} \hi{जा} \hi{रहा} \hi{था।}
\end{quote}

\begin{quote}
\emph{rājmārg par ṭrakoṁ kā ek lambā qāfilā jā rahā thā.}
\end{quote}

\begin{quote}
A long convoy of lorries was moving along the highway.
\end{quote}

Masculine.

\begin{cousinweb}
\textbf{\hi{क़ाफ़िला}} (\emph{qāfilā}) is usually traced to Arabic \emph{qāfila}, \textquotedblleft{}a caravan\textquotedblright{}, which came to Hindi through Persian: once a line of camels and travellers crossing the desert together, now any line of vehicles on the move.
\end{cousinweb}

\begin{grammarlens}[title={What you've built}]
One more word for getting around.
\end{grammarlens}

\subsection*{Your turn}
\begin{itemize}
\raggedright
  \item \emph{Say it:} \emph{qāfilā}
  \item \emph{Say it:} \emph{qāfilā}, once more
  \item \emph{Say it:} \emph{rājmārg par ṭrakoṁ kā ek lambā qāfilā jā rahā thā.} — A long convoy of lorries was moving along the highway.
  \item \emph{Recall:} say the Hindi for a toll, then the Hindi for a checkpoint, then say \emph{qāfilā} again
\end{itemize}

\begin{tcolorbox}[breakable,colback=teal!4,colframe=teal!35!black,title={Before you move on}]
What does \hi{क़ाफ़िला} mean? (\textquotedblleft{}A convoy; a caravan\textquotedblright{}.) Say it once more.
\end{tcolorbox}
\section[\texorpdfstring{\hi{बहाना} (bahānā)}{बहाना (bahānā)}]{\hi{बहाना} (bahānā) — to wash away; to make (water) flow}

\label{lesson:HI-C407-bahana}



\begin{quote}
Before the new one: say the Hindi for a checkpoint, then the Hindi for a convoy.
\end{quote}

\subsection*{You'll want to know: \hi{बहाना}}
\textbf{\hi{बहाना}} — \emph{bahānā} — \textquotedblleft{}to wash away; to make (water) flow\textquotedblright{}.

\begin{quote}
\hi{बाढ़} \hi{ने} \hi{पुराना} \hi{पुल} \hi{बहा} \hi{दिया।}
\end{quote}

\begin{quote}
\emph{bāṛh ne purānā pul bahā diyā.}
\end{quote}

\begin{quote}
The flood washed away the old bridge.
\end{quote}

Past, with \hi{ने} (he or she …): \hi{उसने} \hi{नाली} \hi{में} \hi{पानी} \hi{बहाया}, \emph{usne nālī meṁ pānī bahāyā}.

A different word, the noun \hi{बहाना}, an excuse, is spelled the same way.

\begin{cousinweb}
\textbf{\hi{बहाना}} (\emph{bahānā}) is the causative partner of \hi{बहना}, \textquotedblleft{}to flow\textquotedblright{}, from Sanskrit \emph{vahati}, \textquotedblleft{}carries along\textquotedblright{} (CDIAL 11453): to make water flow, or to let it carry a thing away.
\end{cousinweb}

\begin{grammarlens}[title={What you've built}]
One more word for getting around.
\end{grammarlens}

\subsection*{Your turn}
\begin{itemize}
\raggedright
  \item \emph{Say it:} \emph{bahānā}
  \item \emph{Say it:} \emph{bahānā}, once more
  \item \emph{Say it:} \emph{bāṛh ne purānā pul bahā diyā.} — The flood washed away the old bridge.
  \item \emph{Recall:} say the Hindi for a checkpoint, then the Hindi for a convoy, then say \emph{bahānā} again
\end{itemize}

\begin{tcolorbox}[breakable,colback=teal!4,colframe=teal!35!black,title={Before you move on}]
What does \hi{बहाना} mean? (\textquotedblleft{}To wash away; to make (water) flow\textquotedblright{}.) Say it once more.
\end{tcolorbox}
